\newpage
\begin{center}
    \Huge The Last Supper and Gethsemane \\ 
    \Large Who Was Working Behind the Scenes?
\end{center}

Let's not forget another detail: at the time of the Last Supper, Jesus already knew that one of his Apostles had betrayed him. While he never mentioned Judas by name, he did give some rather broad hints as to his identity, following which the traitor had no other choice but to run off, making use of the loophole his teacher left him:

\begin{smallquote}
    For some thought, because Judas had the money box, that Jesus had said to him, ``Buy those things we need for the feast,'' or that he should give something to the poor. Having received the piece of bread, he then went out immediately. And it was night. (Jn 13:29-30)
\end{smallquote}
Why Jesus never mentioned his traitor by name---whether because he wasn't totally confident in his information and his hints were a way of verifying it, or because he simply didn't want any bloodshed (some of his Apostles were real tough guys, especially Peter)---isn't all that important. A more interesting question is: how did Jesus obtain the information that someone in his inner circle had turned on him?

This question might seem completely asinine. Clearly, Christ knew about it on account of his divine omniscience; besides, how could he acquire that sort of information under the Apostles' noses when they were always by his side? Well, I'm not going to touch on the topic of his ``omniscience,'' but as for Christ's lack of stable contacts outside of his typical social circle, as described by the Evangelists---that, I'm willing to dispute.

First off: why was the Last Supper so secretive?\footnote{In the Russian Orthodox Church, the Last Supper is commonly referred to as the Secret Supper; with this in mind, a more faithful translation of this sentence would be: ``Why was the Secret Supper, well, secret?'' (Z.B.)} What do we know about the house it took place in, the starting point for Christ and the Apostles' subsequent journey to Gethsemane? No need to strain your memory or raid your bookshelf for a copy of the New Testament---it contains absolutely no information on the subject. Which is quite strange; the Evangelists always go out of their way to describe the owners of the houses their teacher stopped at for some reason or another, but as for the place where such a crucial event occurred---not a peep.\footnote{In Christian tradition, the Upper Room---the place where the Last Supper was held---is taken to be located at David's Tomb, which is on the southern slope of Mount Zion (``Jerusalem,'' \emph{Bible Encyclopedia}, vol. 2, pp. 79-82). It is this room that features in, for example, the numerous picturesque paintings depicting this episode. However, by no means can we take this point of view. That David's Tomb was only constructed in the fourth century C.E. is only half the problem; as far as our approach is concerned, it is far more relevant that the Scriptures directly state that the Last Supper took place in a residential building (see below). (K.E.)}

The description of how Christ and the Apostles found this house is worth quoting, in my opinion. On Thursday morning, the Apostles ask their teacher where he is going to have his Passover meal. Then, Christ sends Peter and John out to Jerusalem, supplying them with the following instructions:

\begin{smallquote}
    And He said to them, ``Behold, when you have entered the city, a man will meet you carrying a pitcher of water; follow him into the house which he enters. Then you shall say to the master of the house, `The Teacher says to you, ``Where is the guest room where I may eat the Passover with My disciples?''' Then he will show you a large, furnished upper room; there make ready.'' So they went and found it just as He had said to them, and they prepared the Passover. (Lk 22:10-13)
\end{smallquote}
Quite frankly, this episode doesn't sound like it's from Sacred Scripture so much as \emph{Seventeen Moments of Spring} (or rather, \emph{Aquarium}).\footnote{A memoir by former Soviet military intelligence officer Viktor Suvorov. (Z.B.)} It's pretty clear we're dealing with liaisons, physical and verbal passwords, a safehouse, and a classical technique to detect shadowing (by means of countersurveillance, as carried out by the partner of the ``man with the pitcher''). On that note, it's also clear that Peter and John were sent there specifically to check whether or not the safehouse had been compromised. The owner of the house, in full accordance with the requirements of secrecy (which evidently haven't changed over the past twenty centuries), never once met his visitors---for which reason the Gospels lack any information about him whatsoever.

The Evangelists clearly deemed these facts unworthy of their attention; for us, however, it is quite significant that Jesus clandestinely (or at least discreetly) kept in contact with at least one group of people in Jerusalem with whom the Apostles were not familiar. The members of this group, on the other hand (which included the ``man with the pitcher''), knew the Apostles by sight at the very least.

My interest now stoked by Jesus' clandestine contacts, I started to look over the text of the New Testament in that light, and immediately found a number of promising episodes. Three or four in, however, I had to tell myself, ``Wait a second!''---I felt like I was molding the facts to fit my preconceptions. Any more of that and I'd be just like one of those half-wit conspiracy theorists who can sniff out a Global Judeo-Masonic Plot in a six-point snowflake. But there's still one episode we should look at, since it might tie directly into the events of Holy Week: the Transfiguration of the Lord.

Not long after his third (and last) trip into Jerusalem, Christ, accompanied by Peter, John, and James, ascended a mountain to partake in prayer. A number of events subsequently occurred, from which, as usual, we will refrain from commenting on the outright miracles: the illumination of Jesus' face and clothing, the cloud that hung over those present, and the voice from the heavens. Here is the ``solid residue'' that remains:

\begin{smallquote}
    But Peter and those with him were heavy with sleep; and when they were fully awake, they saw His glory and the two men who stood with Him. Then it happened, as they were parting from Him, that Peter said to Jesus, ``Master, it is good for us to be here; and let us make three tabernacles: one for You, one for Moses, and one for Elijah''---not knowing what he said. (Lk 9:32-33)
\end{smallquote}
However, Jesus never said anything to confirm his disciples' speculation that the men speaking with him were really Moses and Elijah. The most interesting part is what happens afterwards:

\begin{smallquote}
    Now as they came down from the mountain, Jesus commanded them, saying, ``Tell the vision to no one until the Son of Man is risen from the dead.'' (Mt 17:9) \\
    \\
    So they kept this word to themselves, questioning what the rising from the dead meant. (Mk 9:10)
\end{smallquote}
Simply put: three of the Apostles accidentally witnessed a meeting between Jesus and two unknown men, which was clearly not meant for their eyes. Why else would their teacher have demanded that they keep their mouths shut until his dying day?

However, let's return to our general line of inquiry---the events of Holy Thursday. So, how and when could the information on Judas' betrayal have been relayed to Jesus without the Apostles knowing (the information's source being another question entirely)? There are countless ways it could have been transmitted (for example, through a practically undetectable liaison disguised as a beggar on a porch), but in our case, we can make the simplest assumption possible: Jesus found the corresponding message in a previously arranged location in the safehouse---``large, furnished, and prepared'' (Mk 14:15)---following which he accused Judas of betraying him.

Another mystery remains, however. Following the conclusion of the Last Supper, Jesus and the remaining eleven Apostles left Jerusalem and made their flight to Gethsemane---an olive grove at the foot of the Mount of Olives, which lies to the east of the city.\footnote{Some food for thought: the village of Bethany, which was the setting for a number of miraculous events (see above), is located on the same Mount of Olives as the Garden of Gethsemane, albeit the former is on its southeastern slope, while the latter is on its western slope. (K.E.)} Judging from the lengthy conversation they had on the way there, the route was not a short one. It was at Gethsemane, after the Agony in the Garden, that Christ was arrested. Nowhere in the Gospels are there any indications that the Apostles had previously been informed of their teacher's plans---namely, his intended destination after leaving Jerusalem. Judas, on the other hand, who left the Last Supper long before its conclusion, managed to lead the arrest team exactly where they needed to go---the Garden of Gethsemane.

I hope we won't try to ascribe the gift of omniscience to this man, too? Let's think this through logically, then. Now, as for how Judas managed to guess the location, that's more or less clear: apparently, Gethsemane served as a constant refuge for the Apostles over the past few days (``And in the daytime He was teaching in the temple, but at night He went out and stayed on the mountain called Olivet''---Lk 21:37). It should be noted that the Garden is actually home to a crag with a large cave, which presumably served as a place of shelter for the Apostles.\footnote{``Gethsemane,'' \emph{Bible Encyclopedia}, vol. 1, pp. 156-157.} It was ostensibly from there that they headed out to Jerusalem for the Last Supper. Another puzzle: on what grounds was Judas sure that Jesus would return that night to his obviously compromised base? Putting himself in his teacher's place, he ought to have supposed that he would try to evade detection and flee Jerusalem at once, as he had repeatedly done before (for example, Jn 10:39-40). These ``disappearances'' of Jesus are interesting in and of themselves, but they are irrelevant to the discussion at hand.

We've somehow gotten used to assuming that all of Christ's actions (or rather, lack thereof) that night were rooted in his firm intention to ``drink of this cup.'' For actual operational-search measures to succeed, however, detectives need to at least genuinely consider the motivations of those they are pursuing. Do Judas and the Sanhedrin's officers seem like they could penetrate the Savior's thoughts so precisely? The question, as I see it, is a rhetorical one. And on the other hand: suppose Jesus had already made his final arrangements with regard to his own life. However, could a man like him choose to stay alongside his disciples an consciously expose them to a deadly, not to mention completely senseless risk? At the time of his arrest, the Teacher declares, ```If you seek Me, let these go their way' (that the saying might be fulfilled which He spoke, `Of those whom You gave Me I have lost none')'' (Jn 18:8-9). But let's be honest: his own contribution to his disciples remaining alive seems minimal at best. After all, it's simply a ``criminal oversight'' that they weren't slaughtered on the spot during the incident with Malchus!

All these considerations allow us to make the following assumption. That night, Jesus had some task waiting for him at Gethsemane---something so important that, unless he completed it, he could neither leave Jerusalem (even under threat of arrest) or voluntarily give himself up to the high priests; Judas likewise knew of this task. We can assume that in the garden---the place where he had most recently taken shelter---Jesus was to either take something at a predetermined location, or, on the contrary, leave something, or---most likely of all---meet someone. And if there really was someone waiting for him in the garden (recall, for example, the nighttime visit paid to Jesus a year prior by Nicodemus, a member of the Sanhedrin), then it becomes clear why Jesus tried to reach Gethsemane before the detectives, paying no heed to the danger involved.

When was this meeting supposed to take place (whether it occurred in reality is a different question entirely)? I think it was the very moment Jesus retreated alone into the garden to pray---and here's why. Let's compare the Agony in the Garden with the Transfiguration of the Lord (see above). In both cases, Christ breaks off from the disciples so he can pray in solitude. In both cases, he is accompanied by three Apostles---the same three, no less: Peter, John, and James. In both cases, all three ``bodyguards'' strangely fall asleep. Doesn't it sound like there's a few too many coincidences going on here, and the two stories actually describe a single meeting the Teacher took part in, the true nature of which was poorly understood by his Apostles? Especially considering there's also a direct indication that Jesus was not alone in the garden (Lk 22:43); admittedly, the Evangelist believed his teacher was accompanied by an ``angel,'' but that is just pure speculation.

Let's ask ourselves a slightly different question: could the Apostles have possibly made out the person conversing with their teacher further off into the garden (assuming such a meeting actually took place)? I think not, and here's why. Let's recall another event from that night---Peter's threefold denial in Caiaphas' courtyard:

\begin{smallquote}
    Now the servants and officers who had made a fire of coals stood there, for it was cold, and they warmed themselves. And Peter stood with them and warmed himself. (Jn 18:18)
\end{smallquote}
This was in the courtyard of a mansion in the city; you can only imagine how bitterly cold it was that spring night in the Garden of Gethsemane. As a result, it's quite reasonable to suppose that the Apostles also made a fire to keep warm. And in principle, that meant they couldn't make out anything beyond the circle of light.

But of course, could the three ``attendants'' have noticed anything while they were sleeping on the job? For what it's worth, their strange slumber seems understandable and natural when taken in light of the last assumption. You have to think that a few minutes into their vigil, the disciples (side note: didn't the Teacher give these guys the task of keeping watch?) were freezing to the bone and decided---just for a second!---to warm themselves by the fire. Alas, ``seconds'' like these always end the same way---the heat gets you drowsy in a flash, and then it's lights out... So I wouldn't put too much stock in their testimony, either, but even then, they certainly noticed \emph{something}---\emph{before} and \emph{after} their stay by the fire. In any case, the information on the ``angel'' visiting Jesus definitely comes from them---there was simply no one else who could have seen it.